\section{Algoritmos}
\label{sec:algoritmos}

Esta sección presenta en pseudocódigo los cuatro algoritmos centrales del
planificador: el sorteo ponderado, el bucle de despacho del planificador, el
ciclo de un worker (que incluye la cesión y la actualización de boletos
activos de la extensión), y el cálculo del tamaño de bloque cooperativo. El
código real vive en \texttt{src/sync.c}, \texttt{src/scheduler.c} y
\texttt{src/worker.c}.

\begin{algorithm}[htbp]
\caption{Sorteo ponderado (\texttt{sync\_select\_winner})}
\label{alg:sorteo}
\begin{algorithmic}[1]
\REQUIRE tareas $T[1..N]$, generador $\text{rng}$
\STATE \textbf{lock}(mutex)
\STATE $\text{activos} \gets \sum_{i:\, T[i].\text{state}=\text{READY}} T[i].\text{effective\_tickets}$
\IF{$\text{activos} = 0$}
  \STATE \textbf{unlock}(mutex); \textbf{return} $\bot$ \COMMENT{todas terminaron}
\ENDIF
\STATE $\text{boleto} \gets (\text{rng.next}() \bmod \text{activos}) + 1$
\STATE $\text{acum} \gets 0$
\FOR{$i \gets 1$ \TO $N$}
  \IF{$T[i].\text{state} = \text{READY}$}
    \STATE $\text{acum} \gets \text{acum} + T[i].\text{effective\_tickets}$
    \IF{$\text{boleto} \le \text{acum}$}
      \STATE \textbf{unlock}(mutex); \textbf{return} $i$ \COMMENT{ganadora}
    \ENDIF
  \ENDIF
\ENDFOR
\end{algorithmic}
\end{algorithm}

El sesgo del módulo en la línea 6 (cuando \text{activos} no divide exacto el
rango de \texttt{xorshift32}) se acepta y documenta en vez de corregirse con
\textit{rejection sampling}: para los tamaños de \text{activos} reales del
proyecto (decenas a unos pocos miles) el sesgo relativo es del orden de
$10^{-8}$, siete órdenes de magnitud por debajo de la tolerancia del
experimento de proporcionalidad ($\leq 0.02$).

\begin{algorithm}[htbp]
\caption{Bucle del planificador (\texttt{scheduler\_run}, hilo principal)}
\label{alg:bucle}
\begin{algorithmic}[1]
\STATE despachos $\gets 0$
\LOOP
  \IF{$\text{max\_despachos} > 0$ \AND despachos $\ge$ max\_despachos}
    \STATE \texttt{sync\_request\_stop}() \COMMENT{despierta a las READY restantes sin despacharlas}
    \STATE \textbf{break}
  \ENDIF
  \STATE $g \gets$ \texttt{sync\_select\_winner}()
  \IF{$g = \bot$}
    \STATE \textbf{break} \COMMENT{ninguna tarea READY: todas terminaron}
  \ENDIF
  \STATE \texttt{sync\_dispatch}($g$) \COMMENT{state[$g$]=RUNNING, signal dirigida a $g$}
  \STATE \texttt{sync\_wait\_for\_event}() \COMMENT{bloquea hasta que $g$ ceda o termine}
  \STATE despachos $\gets$ despachos $+ 1$
  \STATE registrar evento de despacho (log)
\ENDLOOP
\STATE \texttt{worker\_pool\_join}()
\end{algorithmic}
\end{algorithm}

\begin{algorithm}[htbp]
\caption{Ciclo del worker de la tarea $t$ (\texttt{worker\_thread\_main})}
\label{alg:worker}
\begin{algorithmic}[1]
\LOOP
  \STATE \texttt{sync\_wait\_for\_dispatch}($t$) \COMMENT{bloquea hasta RUNNING o stop}
  \IF{el planificador solicitó parada antes de despachar a $t$}
    \STATE \textbf{return} \COMMENT{$t$ queda READY, sin ejecutar nada}
  \ENDIF
  \STATE $\text{bloque} \gets$ tamaño de bloque del modo (Algoritmo~\ref{alg:cooperativo} en modo cooperativo, o quantum $Q$ fijo)
  \STATE $u \gets$ \textsc{ActualizarBoletos}($t$, bloque) \COMMENT{Sección~\ref{sec:extension}; sin \texttt{--yield-config}, $u=\min(\text{bloque}, \text{restante})$}
  \STATE \texttt{workload\_run\_units}($t$, $u$) \COMMENT{fuera del mutex}
  \IF{$t.\text{completado} = t.\text{trabajo\_total}$}
    \STATE $\text{siguiente} \gets \text{FINISHED}$
  \ELSE
    \STATE $\text{siguiente} \gets \text{READY}$
  \ENDIF
  \STATE \texttt{sync\_finish\_turn}($t$, siguiente) \COMMENT{signal a \texttt{cond\_scheduler}}
  \IF{$\text{siguiente} = \text{FINISHED}$}
    \STATE \textbf{return}
  \ENDIF
\ENDLOOP
\end{algorithmic}
\end{algorithm}

\begin{algorithm}[htbp]
\caption{Tamaño de bloque cooperativo (\texttt{cooperative\_slice\_size})}
\label{alg:cooperativo}
\begin{algorithmic}[1]
\REQUIRE $\text{trabajo\_total} \ge 1$, $1 \le P \le 100$
\STATE \textbf{return} $\max\left(1,\ \left\lceil \dfrac{\text{trabajo\_total} \times P}{100} \right\rceil\right)$
\end{algorithmic}
\end{algorithm}

El bloque cooperativo se calcula una sola vez sobre el \texttt{work\_units}
\emph{total} de la tarea (no sobre el trabajo restante) y es constante
durante toda su vida; en cada activación se topa por
$\min(\text{bloque}, \text{restante})$, igual que el quantum fijo. Ambos
modos comparten exactamente la misma función de ejecución
(\texttt{workload\_run\_units}) y el mismo bucle del worker --- la única
diferencia mecánica entre cooperativo y quantum es cómo se calcula el
tamaño del bloque, lo que explica por qué el Caso~6 (Sección~\ref{sec:pruebas})
puede exigir resultados idénticos entre ambos modos con solo distinto costo
de coordinación. La división se implementa con aritmética entera
($\lceil a/b \rceil = \lfloor (a+b-1)/b \rfloor$) en vez de \texttt{double},
para preservar el determinismo bit a bit exigido por el enunciado.
